\section*{الفصل \ref{ch:g12:cells-and-electrolysis} --- الخلايا الكهروكيميائية والتحليل الكهربائي}

\begin{solution}{exo:g12:cells-and-electrolysis:1}
قطب الخارصين هو \omterm{def:g12:cells-and-electrolysis:anode}{المصعد}، \ce{Zn -> Zn^{2+} + 2e-}؛ وقطب النحاس
هو \omterm{def:g12:cells-and-electrolysis:anode}{المهبط}، \ce{Cu^{2+} + 2e- -> Cu}.
\end{solution}

\begin{solution}{exo:g12:cells-and-electrolysis:2}
% ledger: const:F
$Q = 0.20 \times 1800 = \qty{360}{C}$؛ $n(e^-) = 360/96485 =
\qty{3.7e-3}{mol}$.
\end{solution}

\begin{solution}{exo:g12:cells-and-electrolysis:3}
يغلق الدارة ويحفظ كل \omterm{def:g2:mixing-and-dissolving:dissolve}{محلول} متعادلًا كهربائيًا، إذ تنتقل
أيوناته إلى الحجرتين. ودونه تكون الدارة مفتوحة: فلا
يسري تيار، ويقرأ مقياس الفولت صفرًا ويبقى المصباح مطفأً.
\end{solution}

\begin{solution}{exo:g12:cells-and-electrolysis:4}
$2.5 \times 3600 = \qty{9000}{C}$.
\end{solution}

\begin{solution}{exo:g12:cells-and-electrolysis:5}
لا: فالتفاعل يجري في الاتجاه المعاكس للاتجاه التلقائي.
والمولّد هو الذي يمد بالطاقة.
\end{solution}

\begin{solution}{exo:g12:cells-and-electrolysis:6}
\omterm{def:g12:cells-and-electrolysis:anode}{المصعد} (النحاس، $-$): \ce{Cu -> Cu^{2+} + 2e-}. \omterm{def:g12:cells-and-electrolysis:anode}{المهبط} (الفضة، $+$):
\ce{Ag+ + e- -> Ag}. والتفاعل الإجمالي: \ce{Cu + 2Ag+ -> Cu^{2+} + 2Ag}.
\end{solution}

\begin{solution}{exo:g12:cells-and-electrolysis:7}
ثلاث سنوات تساوي $3 \times 365 \times 24 = \qty{26280}{h}$؛
$0.010 \times 26280 = \qty{263}{mA.h}$.
\end{solution}

\begin{solution}{exo:g12:cells-and-electrolysis:8}
% ledger: const:F, aw:Ag
$Q = 0.50 \times 2400 = \qty{1200}{C}$؛ $n(e^-) = n(\ce{Ag}) =
1200/96485 = \qty{0.0124}{mol}$؛ $m = 0.0124 \times 107.9 =
\qty{1.34}{g}$.
\end{solution}

\begin{solution}{exo:g12:cells-and-electrolysis:9}
يسري التيار في السلك من النحاس ($+$) إلى الخارصين ($-$)،
عكس \omterm{def:g9:inside-the-atom:nucleus}{الإلكترونات}. وتتحرك \omterm{def:g9:ions:ion}{أيونات} \ce{K+} نحو حجرة
النحاس، التي تفقد \ce{Cu^{2+}}؛ \omterm{def:g9:ions:ion}{وأيونات} \ce{NO3-} نحو حجرة
الخارصين، التي تكسب \ce{Zn^{2+}}.
\end{solution}

\begin{solution}{exo:g12:cells-and-electrolysis:10}
في كلتا الحالتين \omterm{def:g12:cells-and-electrolysis:anode}{المصعد} هو حيث تجري \omterm{def:g11:redox:oxidation}{الأكسدة}. ففي الخلية يكون \omterm{def:g12:cells-and-electrolysis:anode}{المصعد}
$-$، والتفاعل تلقائيًا، والطاقة الكيميائية تتحول إلى
طاقة كهربائية؛ وفي \omterm{def:g12:cells-and-electrolysis:electrolysis}{التحليل الكهربائي} يوصل \omterm{def:g12:cells-and-electrolysis:anode}{المصعد} بالقطب $+$
للمولّد، والتفاعل مفروضًا، والطاقة الكهربائية تتحول إلى
طاقة كيميائية.
\end{solution}

\begin{solution}{exo:g12:cells-and-electrolysis:11}
\qty{12}{mL}: فالمعادلة \ce{2H2O -> 2H2 + O2} تعطي جزيئين من
ثنائي الهيدروجين مقابل \omterm{def:g7:atoms-and-molecules:molecule}{جزيء} من ثنائي الأكسجين، والكميات المتساوية من الغاز تشغل
أحجامًا متساوية.
\end{solution}

\begin{solution}{exo:g12:cells-and-electrolysis:12}
% ledger: aw:Zn, const:F
$n(\ce{Zn}) = 5.0/65.4 = \qty{0.076}{mol}$؛
$n(\ce{Cu^{2+}}) = 0.100 \times 0.50 = \qty{0.050}{mol}$؛ ويتفاعلان
واحدًا لواحد، فتنفد \omterm{def:g9:ions:ion}{أيونات} النحاس أولًا. $n(e^-) = 2 \times 0.050
= \qty{0.100}{mol}$، $Q = \qty{9.65e3}{C} = \qty{2.68}{A.h}$. وعند
\qty{50}{mA}: $2.68/0.050 = \qty{54}{h}$.
\end{solution}

\begin{solution}{exo:g12:cells-and-electrolysis:13}
$Q = 2.0 \times 3600 = \qty{7200}{C}$؛ $E = 7200 \times 1.5 =
\qty{1.08e4}{J}$، أي $1.5 \times 2.0 = \qty{3.0}{W.h}$.
\end{solution}

\begin{solution}{exo:g12:cells-and-electrolysis:14}
% ledger: const:F
$Q = \qty{7200}{C}$، $n(e^-) = \qty{0.0746}{mol}$. إلكترونان لكل
\ce{H2}: $n(\ce{H2}) = \qty{0.0373}{mol}$، $V = \qty{0.90}{L}$. وأربعة لكل
\ce{O2}: $n(\ce{O2}) = \qty{0.0187}{mol}$، $V = \qty{0.45}{L}$.
\end{solution}

\begin{solution}{exo:g12:cells-and-electrolysis:15}
النحاس الذائب عند \omterm{def:g12:cells-and-electrolysis:anode}{المصعد} يساوي النحاس المترسب، لأن \omterm{def:g9:inside-the-atom:nucleus}{الإلكترونات}
نفسها تمر عند القطبين: $2.84/3.00 = \qty{94.7}{\percent}$ من النحاس.
\end{solution}

\begin{solution}{pb:g12:cells-and-electrolysis:1}
% ledger: const:F, const:NA, aw:Li, dch:CH4
\textbf{1.} \omterm{def:g12:cells-and-electrolysis:anode}{المصعد}: فالليثيوم يتأكسد عنده.

\textbf{2.} قطب الغرافيت، أي \omterm{def:g12:cells-and-electrolysis:anode}{المصعد}.

\textbf{3.} تنتقل \omterm{def:g9:inside-the-atom:nucleus}{الإلكترونات} عبر الهاتف من قطب الغرافيت إلى
قطب الأكسيد؛ وتعبر \omterm{def:g9:ions:ion}{أيونات} الليثيوم الإلكتروليت في الاتجاه
نفسه، من الغرافيت إلى الأكسيد.

\textbf{4.} يمكن دفع تفاعلها في الاتجاه المعاكس بشاحن.

\textbf{5.} \qty{4.000}{A.h}، أي $4.000 \times 3600 =
\qty{14400}{C}$.

\textbf{6.} $14400/96485 = \qty{0.1492}{mol}$.

\textbf{7.} $0.1492 \times \num{6.022e23} = \num{8.99e22}$ \omterm{def:g9:inside-the-atom:nucleus}{إلكترون}.

\textbf{8.} $4.000/0.25 = \qty{16}{h}$.

\textbf{9.} $0.20 \times 14400 = \qty{2880}{C}$.

\textbf{10.} $E = 14400 \times 3.7 = \qty{5.3e4}{J}$، أي نحو \qty{53}{kJ}.

\textbf{11.} $53280/3600 = \qty{14.8}{W.h}$.

\textbf{12.} حرق \qty{1}{g} من الميثان يطلق طاقة أكبر قليلًا
(\qty{55.7}{kJ}) مما تخزنه بطارية هاتف ممتلئة.

\textbf{13.} $1000/14.8 = 67.6$: أي 67 شحنة كاملة.

\textbf{14.} \ce{Li+ + e- -> Li}: \omterm{def:g11:redox:oxidation}{اختزال}.

\textbf{15.} يفرض الشاحن التفاعل في الاتجاه المعاكس
للاتجاه التلقائي.

\textbf{16.} $14400/2.0 = \qty{7200}{s}$، أي \qty{2.0}{h}.

\textbf{17.} \omterm{def:g9:ions:ion}{أيون} ليثيوم واحد لكل \omterm{def:g9:inside-the-atom:nucleus}{إلكترون}: \qty{0.1492}{mol}.

\textbf{18.} $0.1492 \times 6.9 = \qty{1.03}{g}$.

\textbf{19.} يتنقل نحو \qty{1.03}{g} من الليثيوم بين
القطبين في شحنة كاملة واحدة.
\end{solution}
